Okay, let's tackle this problem step by step. The user wants me to analyze an algorithm that uses minibatch gradients with a fixed step size for non-convex and smooth functions. They need a detailed LaTeX document following specific sections and structures based on the examples provided. 

First, I need to identify the algorithm. The description mentions using randomly sampled minibatches and taking a step in the negative gradient direction with a fixed step size. That sounds like Mini-batch Gradient Descent (MGD). So the algorithm name is likely Mini-batch Gradient Descent.

Next, the algorithm details. The update rule would involve taking a minibatch gradient. Let's formalize that. The parameters are updated as w_{t+1} = w_t - η * g_t, where g_t is the average gradient over the minibatch. The learning rate η is fixed. The pseudocode should include initializing w, looping over iterations, sampling minibatches, computing gradients, updating weights, and checking termination conditions.

For the dry run example with f(w) = (w - 3)^2, starting at w0=0 and η=0.1. Let's compute each iteration. The gradient of f is 2(w - 3). For each minibatch, assuming the entire dataset is used each time (which is a common simplification for such examples), the gradient at w=0 is 2*(0-3) = -6. So the update would be w1 = 0 - 0.1*(-6) = 0.6. Then compute f(w1)= (0.6 -3)^2 = 6.76. Next iteration, gradient at 0.6 is 2*(0.6-3)= -4.4, update to 0.6 + 0.44 = 1.04. f(w)= (1.04-3)^2=4.3024. Third iteration: gradient -3.72, update to 1.04 + 0.372=1.412. f(w)= (1.412-3)^2≈2.673. So the example should show these steps.

Moving to convergence theory. Assumptions for non-convex and smooth functions. The function f should be L-smooth, meaning the gradient is Lipschitz continuous. Also, the stochastic gradients should be unbiased and have bounded variance. The convergence theorem would state that the expected norm of the gradient decreases by a factor of (1 - η L)^2 each step, assuming η ≤ 1/L. However, for the rate, in the non-convex case, we might get O(1/√T) average gradient norm, similar to SGD but with different constants.

In the proof, I need to start with the smoothness assumption. For each step, expand the distance squared to the optimal point, but in non-convex case, we can't assume a single minimum. So instead, we consider the gradient at the current point. Using the smoothness, we can bound the difference in function values. Then take expectations over the stochastic gradients. The key steps would involve expanding the update, applying the smoothness condition, and taking expectations. Each step needs to be annotated with line numbers for clarity.

Also, need to mention hyperparameter sensitivity, like how the learning rate η affects stability. If η is too large, the algorithm diverges; if too small, convergence is slow. Compare to SGD, which has similar convergence rates but with adaptive stepsizes sometimes.

Wait, the user mentioned non-convex and smooth functions. So the convergence theorem might be about the average gradient norm rather than function value. For example, after T iterations, the average of ||∇f(w_t)||² is O(1/√T), similar to SGD but with possibly different constants.

Putting this all together, ensuring each section is covered as per the structure. The LaTeX code should be complete, with sections, equations, and proper formatting. Need to make sure all the required parts are included, like the dry run example with exact numbers, and the proof with step annotations.
</think>

\documentclass{article}
\usepackage{amsmath,amssymb,amsthm}

\section*{Mini-batch Gradient Descent (MGD)}
\section*{Mathematical Formulation}
\begin{itemize}
    \item Parameters: $w \in \mathbb{R}^d$, objective function $f(w)$
    \item Minibatch size: $b \in \mathbb{N}$, batch gradient: $g_t = \frac{1}{b}\sum_{i \in \mathcal{B}_t} \nabla f(w_t; x_i)$
    \item Update rule: $w_{t+1} = w_t - \eta g_t$ where $\eta > 0$ is fixed learning rate
    \item Hyperparameters: $\eta$, $b$, termination after $T$ iterations
\end{itemize}

\section*{Pseudocode}
\begin{verbatim}
Initialize w_0 in domain
for t = 0 to T-1:
    Sample minibatch ᶠB_t of size b
    Compute g_t = (1/b) * ∇_w f(w_t; ᶠB_t)
    w_{t+1} = w_t - η * g_t
return w_T
\end{verbatim}

\section*{Dry Run Example}
Let $f(w) = (w-3)^2$, initial $w_0=0$, $\eta=0.1$, $b=